% Generated draft: requires mathematical and source review.
\documentclass[11pt]{article}
\usepackage[margin=25mm]{geometry}
\usepackage{fontspec}
\setmainfont{Latin Modern Roman}
\usepackage{amsmath,amssymb,mathtools,unicode-math}
\setmathfont{Latin Modern Math}
\usepackage{enumitem,needspace,xurl,hyperref}
\setlength{\parindent}{0pt}
\setlength{\parskip}{5pt}
\setlength{\emergencystretch}{3em}
\newcommand{\N}{\mathbb{N}}
\newcommand{\Z}{\mathbb{Z}}
\newcommand{\Q}{\mathbb{Q}}
\newcommand{\R}{\mathbb{R}}
\newcommand{\C}{\mathbb{C}}
\newcommand{\F}{\mathbb{F}}
\newcommand{\A}{\mathbb{A}}
\newcommand{\PP}{\mathbb{P}}
\newcommand{\E}{\mathbb{E}}
\newcommand{\eps}{\varepsilon}
\newcommand{\dd}{\,\mathrm{d}}
\begin{document}
\section*{UnsolvedMath ID: 20003094}
\textbf{Source problem code:} catalogue entry
\subsection{Definitions and mathematical statement}
Let $S$ be a finite-type hyperbolic surface in the intended Thurston-metric setting and $X,Y\in\mathcal T(S)$. The source asks to identify essential simple closed curves which become short on a preferred Thurston geodesic $G_{X,Y}$. It attributes to F. Kassel and K. Rafi a path linear in shear/cataclysm coordinates for a canonical lamination $\lambda(X,Y)$, with all curve lengths convex after the specified reparametrization. The exact preferred-path construction and numerical meaning of short are not supplied. For a chosen threshold $\epsilon>0$ write $m_\alpha=\min_{Z\in G_{X,Y}}\ell_\alpha(Z)$; the question is to detect $m_\alpha<\epsilon$ from endpoint data.

The supplied partial result assumes endpoints thick with a fixed lower systole bound, and a curve interacting with the maximally stretched lamination in the sense of the cited DLRT theorem. Let $d=d_\alpha(X,Y)$ be its endpoint annular twisting quantity, defined using that theorem's markings/conventions. Its coarse estimate is written with constants $K\geq1$, $C\geq0$ as
\[
K^{-1}F(m_\alpha)-C\leq d\leq K F(m_\alpha)+C,\qquad
F(m)=m^{-1}\max\{1,\log(m^{-1})\}.
\]
The constants depend on the theorem's surface/thickness and normalization data and are not numerically supplied. Since $F$ is decreasing, the exact threshold consequences are
\[
d<K^{-1}F(\epsilon)-C\ \Longrightarrow\ m_\alpha>\epsilon,
\qquad
d>KF(\epsilon)+C\ \Longrightarrow\ m_\alpha<\epsilon.
\]
These leave an intermediate undecided range and do not constitute a necessary-and-sufficient threshold test with specified effective constants.

For sufficiently large $d$, so $(d-C)/K\geq e$, invert $F$ using the nonnegative real Lambert function $W_0$, defined by $W_0(x)e^{W_0(x)}=x$. The aid obtains
\[
\frac{W_0(K(d+C))}{K(d+C)}\leq m_\alpha\leq
\frac{W_0((d-C)/K)}{(d-C)/K}.
\]
Thus the large-twisting scale is comparable to $(\log d)/d$, with constants under the same hypotheses, not an exact asymptotic constant.

For the once-punctured torus $S_{1,1}$ the aid reports that the interacting-curve theorem covers every essential curve in its nontrivial path setting. With the preferred common shear parameter, the short-time set $I_\alpha=\{t:\ell_\alpha(G(t))<\epsilon\}$ is an interval by convexity. For a collar threshold, for example $0<\epsilon<2\operatorname{arsinh}(1)$, two distinct essential curves cannot both be that short, since they intersect and satisfy the collar inequality. Hence the $I_\alpha$ are pairwise disjoint. The further reported order by large-coefficient Farey pivots uses the torus correspondence, not convexity alone; its exact markings, pivots, and constants are absent. Thin endpoints and higher-complexity curves outside the interaction hypothesis remain beyond the supplied result.
\subsection{Short English statement}
For thick endpoints and curves covered by the supplied twisting theorem, endpoint annular twisting gives two one-sided shortness certificates and explicit Lambert-function length bounds. On the punctured torus the preferred-path convexity and collar lemma make short-curve times disjoint intervals; their reported Farey order requires additional torus data. The result does not cover arbitrary thin endpoints or noninteracting curves.
\subsection{Sources}
\begin{itemize}
\item Local registry: permanent catalogue entry, exact source code catalogue entry; \texttt{aim-topology-notes.json}, \texttt{notes[181]}.
\item AIM source \emph{Problems on Thurston metric}, dated 21 April 2014, page 2, Problem 2.4 attributed to K. Rafi; workshop catalog title \emph{Lipschitz metric on Teichmueller space}, \url{https://aimath.org/pastworkshops/lipschitzteichproblems.pdf}. The original text attributes the preferred shear-linear convex path to F. Kassel and K. Rafi.
\item Supplied attribution: DLRT coarse twisting/minimum-length theorem and punctured-torus correspondence, without expanded names, title, theorem number, or complete authoritative citation.
\item Supplied research attempt 1, registry updated 14 August 2026: partial result, low novelty confidence, machine-generated and not peer reviewed. It supplies threshold inversion and the torus synthesis. Verification reports direct inspection of source page 2 in that inherited run; no new inspection is claimed. Original remarks and literature field are empty.
\end{itemize}
\end{document}
